\documentclass[11pt]{article}
\input{style-pc}

\begin{document}

\entetesujet{Sujet bac 2012 -- Physique-Chimie -- Série C}

% ============================ CHIMIE ============================
\matiere{CHIMIE}{8 points}

\exo{1}{(Noyau atomique -- Radioactivité)}
Le polonium $\ch{^{210}_{84}Po}$, noyau instable, se désintègre suivant le mode $\alpha$ en donnant le noyau de plomb (Pb) dans son état fondamental.

\begin{enumerate}
  \item Calculer, en MeV, l'énergie de liaison par nucléon du noyau de polonium.
  \item Écrire l'équation-bilan de la réaction de désintégration d'un noyau de polonium en précisant les lois de conservation utilisées.
  \item Calculer en MeV, l'énergie libérée lors de cette désintégration.
  \item La période du nucléide $\ch{^{210}_{84}Po}$ est $T = 138$ jours. Un échantillon de polonium 210 a une masse initiale $m_0 = 20$ g.
  \begin{enumerate}[label=\alph*.]
    \item Calculer le nombre $N_0$ de noyaux de polonium 210 correspondant.
    \item Calculer la masse de polonium disparu au bout de 414 jours.
  \end{enumerate}
\end{enumerate}
\noindent On donne : $m(\ch{Po}) = 209{,}9369\,\mu$ ; $m(\ch{Pb}) = 205{,}9296\,\mu$ ; $m(\ch{He}) = 4{,}0015\,\mu$. $1\,\mu = 931{,}5$ MeV/C$^2$ ; $\aleph = 6{,}02\cdot10^{23}$ mol$^{-1}$ ; masse proton : $m_p = 1{,}00727\,\mu$ ; $c = 3\cdot10^8$ m$\cdot$s$^{-1}$ ; masse neutron : $m_n = 1{,}00866\,\mu$ ; $M(\ch{Po}) = 210$ g$\cdot$mol$^{-1}$.

\exo{2}{(Couple Acide/Base dans l'eau)}
La méthanamine $\ch{CH_3NH_2}$ est une base dont l'acide conjugué est l'ion méthanammonium $\ch{CH_3NH_3^+}$.

\begin{enumerate}
  \item Écrire l'équation de la réaction de la méthanamine sur l'eau.
  \item On prépare un mélange contenant $20$ cm$^3$ d'une solution de méthanamine de concentration $C_1 = 0{,}1$ mol$\cdot$L$^{-1}$ et $10$ cm$^3$ d'une solution de chlorure de méthanammonium $(\ch{CH_3NH_3^+} + \ch{Cl^-})$ de concentration $C_2 = 0{,}2$ mol$\cdot$L$^{-1}$. Le pH de la solution obtenue est $10{,}6$.
  \begin{enumerate}[label=\alph*.]
    \item Recenser les espèces chimiques présentes dans la solution.
    \item Calculer la concentration molaire de chaque espèce.
    \item En déduire le pKa du couple $\ch{CH_3NH_3^+/CH_3NH_2}$.
  \end{enumerate}
  \item Le pKa du couple $\ch{NH_4^+/NH_3}$ vaut $9{,}2$. Laquelle des deux bases $\ch{CH_3NH_2}$ et $\ch{NH_3}$ est la plus forte ? Justifier.
\end{enumerate}

% ============================ PHYSIQUE ============================
\matiere{PHYSIQUE}{12 points}

\exo{1}{(Dynamique)}
Un camion dont la masse totale a pour valeur $M = 7$ tonnes démarre sur une route rectiligne et horizontale. Il atteint une vitesse de $60$ km/h en $4$ min et continue ensuite à une vitesse constante.

Dans cette question et toutes celles qui suivent, on admettra que l'ensemble des forces de frottement et de résistance de l'air est équivalent à une force unique opposée à la vitesse, d'intensité constante $f = 500$ N.

\begin{enumerate}
  \item Calculer l'intensité de la force de traction développée par le moteur :
  \begin{enumerate}[label=\alph*.]
    \item Au cours du démarrage (le mouvement étant alors supposé rectiligne et uniformément accéléré).
    \item Quand le mouvement est rectiligne et uniforme.
  \end{enumerate}
  \item Pour arrêter le camion, le chauffeur débraie, supprimant ainsi la liaison entre le moteur et les roues motrices pour annuler la force de traction, et en même temps il serre les freins. Le camion, qui roulait à la vitesse de $60$ km/h, s'arrête sur un parcours de $200$ m. Calculer :
  \begin{enumerate}[label=\alph*.]
    \item L'intensité de la force de freinage.
    \item Le temps mis par le camion pour s'arrêter.
  \end{enumerate}
\end{enumerate}

\exo{2}{(Propagation des ondes)}
L'extrémité $O$ d'une corde vibrante est animée d'un mouvement sinusoïdal dont l'équation est : $y_0(t) = 2\cdot10^{-2}\sin 200\pi t$.

\begin{enumerate}
  \item En déduire la fréquence et l'amplitude du mouvement.
  \item La distance qui sépare deux points successifs qui vibrent en opposition de phase est $d = 20$ cm. Calculer :
  \begin{enumerate}[label=\alph*.]
    \item La longueur d'onde.
    \item La vitesse de propagation des ondes.
  \end{enumerate}
  \item Soit $M$ le premier point qui vibre en quadrature de phase avec la source $O$.
  \begin{enumerate}[label=\alph*.]
    \item Déterminer la distance $x$ par rapport à la source $O$.
    \item Établir l'équation horaire du point $M$.
    \item Représenter graphiquement dans un même système d'axes $y_0(t)$ et $y_M(t)$.
  \end{enumerate}
\end{enumerate}

\exo{3}{(Courant alternatif)}
Un circuit électrique est constitué d'un conducteur ohmique de résistance $R$ et d'une bobine d'inductance $L$ et de résistance négligeable.

\begin{enumerate}
  \item On alimente ce circuit sous une tension continue $U_1 = 6$ V, l'intensité du courant est $I_1 = 0{,}2$ A. Déterminer la résistance $R$ et la puissance électrique consommée.
  \item Le circuit est ensuite alimenté sous une tension alternative de valeur efficace $U_2 = 6$ V et de fréquence $N = 50$ Hz. L'intensité efficace du courant est $I_2 = 0{,}1$ A. Calculer :
  \begin{enumerate}[label=\alph*.]
    \item La puissance électrique du circuit.
    \item Le facteur de puissance du circuit.
    \item L'inductance $L$ de la bobine.
  \end{enumerate}
  \item Un condensateur associé en série ramène le facteur de puissance du circuit à $0{,}8$. En admettant que le circuit est capacitif, calculer :
  \begin{enumerate}[label=\alph*.]
    \item L'impédance du circuit.
    \item Sa réactance $X$.
    \item La valeur de la capacité $C$ du condensateur.
  \end{enumerate}
\end{enumerate}

\end{document}
